% =============================================================
% บทที่ 3 - วิธีการดำเนินงาน
% เนื้อหาแปลง/เรียบเรียงจาก Progress/Proposal ProjectNew.docx ฉบับปรับปรุงล่าสุด (หัวข้อ 5)
% แผนการดำเนินงาน (Gantt) อัปเดตตามหัวข้อ 9 ของ proposal ฉบับใหม่ (5 เดือน ปี 2569)
% ฐานข้อมูล (PostgreSQL + pgvector / Supabase) อัปเดตตาม proposal ฉบับใหม่ แทนที่ MongoDB+Chroma เดิม
% =============================================================

% -----------------------------------------------------------
\section{ขั้นตอนการดำเนินงาน}
% -----------------------------------------------------------

การดำเนินงานวิจัยนี้แบ่งออกเป็น 5 ส่วนหลัก เพื่อให้การพัฒนาแอปพลิเคชันเป็นไปอย่างมีระบบและครอบคลุมทั้งด้านเทคโนโลยี การจัดการข้อมูล การออกแบบสถาปัตยกรรมระบบ และการประเมินผลลัพธ์ โดยเนื้อหาในส่วนนี้อ้างอิงจากการออกแบบและพัฒนาต้นแบบระบบ (Prototype) ที่ดำเนินการจริงในงานวิจัย ได้แก่ เว็บแอปพลิเคชันแนะนำการท่องเที่ยวจังหวัดขอนแก่นที่ประกอบด้วยระบบ แชตบอต ให้คำแนะนำและระบบวางแผนการเดินทางที่ขับเคลื่อนด้วยโมเดลภาษาขนาดใหญ่ (LLM) ร่วมกับเทคนิค RAG

\subsection{ศึกษาปัญหาและรวบรวมความต้องการ}

\subsubsection{การศึกษาวรรณกรรมและทฤษฎีที่เกี่ยวข้อง}

ขั้นตอนนี้มีวัตถุประสงค์เพื่อรวบรวมและวิเคราะห์แนวคิด เทคนิค และเครื่องมือที่ทันสมัยในการพัฒนาระบบแนะนำการท่องเที่ยว (Tourism Recommender Systems: TRS) และระบบอัจฉริยะ เพื่อใช้เป็นกรอบแนวคิดในการออกแบบระบบ

\begin{enumerate}
\item การศึกษาระบบแนะนำการท่องเที่ยวและการวางแผนการเดินทาง: ศึกษาแนวคิดพื้นฐานของระบบ TRS และวิวัฒนาการไปสู่ระบบที่พิจารณาหลายภาคส่วน (Multistakeholder Recommendation) โดยเฉพาะแนวโน้มการบูรณาการปัจจัยด้านความยั่งยืน (Sustainability) เข้าไปในการวางแผนการเดินทาง ซึ่งครอบคลุมทั้งมิติเศรษฐกิจ สังคม และสิ่งแวดล้อม แนวคิดนี้ถูกนำมาใช้เป็นกรอบอ้างอิงในการออกแบบเกณฑ์การให้คำแนะนำของระบบ กล่าวคือ การจัดลำดับและคัดเลือกสถานที่ในแผนการเดินทางของงานวิจัยนี้พิจารณาปัจจัยด้านงบประมาณ ระยะเวลา ความสนใจของผู้ใช้ และระยะทาง/เวลาการเดินทางระหว่างสถานที่ร่วมกัน โดยยังไม่ได้นำปัจจัยด้านผลกระทบสิ่งแวดล้อมเชิงปริมาณ (เช่น Carbon Footprint) มาใช้ในการคำนวณโดยตรงในต้นแบบปัจจุบัน
\item การศึกษาเทคโนโลยีปัญญาประดิษฐ์สำหรับการประมวลผลภาษาธรรมชาติ: ศึกษานวัตกรรมที่ใช้โมเดลภาษาขนาดใหญ่ (LLMs) และการประมวลผลภาษาธรรมชาติ (NLP) สำหรับการสร้างคำแนะนำและการตอบโต้กับผู้ใช้งาน โดยเน้นที่สถาปัตยกรรม RAG ซึ่งเป็นแนวทางที่มีประสิทธิภาพในการจัดการข้อมูลการท่องเที่ยวที่มีการเปลี่ยนแปลงอยู่เสมอ (Dynamic Data) โดยไม่จำเป็นต้องปรับจูนโมเดล LLM (Fine-tuning) ซ้ำทุกครั้งที่ข้อมูลเปลี่ยนแปลง แนวคิดนี้ถูกนำไปประยุกต์ใช้จริงในทั้งส่วนของระบบ แชตบอต และระบบวางแผนการเดินทางของงานวิจัยนี้
\needspace{6\baselineskip}
\item การศึกษาสถาปัตยกรรมระบบสำหรับแอปพลิเคชัน AI: ศึกษาแนวทางการออกแบบระบบแบบแยกส่วนบริการ (Service-Oriented / Microservices Architecture) ซึ่งช่วยให้แต่ละองค์ประกอบของระบบสามารถพัฒนา ทดสอบ และปรับขนาด (Scale) ได้อย่างเป็นอิสระต่อกัน อันเหมาะสมกับแอปพลิเคชัน AI ที่มีภาระการประมวลผลแตกต่างกันในแต่ละส่วน (เช่น การประมวลผลภาษาธรรมชาติที่ใช้ทรัพยากรสูงกว่าการให้บริการข้อมูลทั่วไป) แนวคิดนี้ถูกนำมาใช้เป็นหลักในการออกแบบระบบจริงของงานวิจัย
\end{enumerate}

\subsubsection{การเก็บรวบรวมและเตรียมข้อมูล}

การเตรียมข้อมูลถือเป็นรากฐานสำคัญในการสร้างระบบอัจฉริยะที่แม่นยำและน่าเชื่อถือ

\begin{enumerate}
\item การรวบรวมข้อมูลแหล่งท่องเที่ยวและกิจกรรม: ข้อมูลหลักของระบบรวบรวมจาก Google Places API (New) ครอบคลุมข้อมูลสถานที่ท่องเที่ยว โรงแรม/ที่พัก และร้านอาหารในพื้นที่จังหวัดขอนแก่น โดยข้อมูลที่ดึงมาประกอบด้วยชื่อสถานที่ ประเภท/แท็ก พิกัดภูมิศาสตร์ คะแนนรีวิว ระดับราคา เวลาทำการ (รวมช่วงเวลาเปิด--ปิดแบบละเอียดรายวัน) และสถานะทางธุรกิจ (Business Status) การนำเข้าข้อมูลตั้งต้นดำเนินการผ่านสคริปต์ดึงและนำเข้าข้อมูล ส่วนการปรับปรุงภายหลังดำเนินการผ่านหน้าจอผู้ดูแลระบบ (Admin Dashboard) ที่พัฒนาขึ้นเป็นส่วนหนึ่งของระบบ ซึ่งผู้ดูแลระบบสามารถเพิ่มหรือแก้ไขข้อมูลสถานที่ กิจกรรม และฐานความรู้ และนำเข้าสถานที่ใหม่จาก Google Places ได้ 2 วิธี คือ ค้นจากชื่อสถานที่ทีละแห่ง หรือระบุชื่อพื้นที่ (เช่น ตำบล อำเภอ หรือย่าน) พร้อมรัศมีและหมวดหมู่ แล้วเลือกนำเข้าหลายแห่งพร้อมกันเป็นชุดงาน โดยให้ AI จัดหมวดหมู่และร่างคำอธิบายให้สถานที่ใหม่ได้ ส่วนข้อมูลที่มีอยู่แล้ว ผู้ดูแลระบบสั่งซิงก์ได้เอง และตั้งให้ระบบซิงก์อัตโนมัติตามรอบเวลาได้ โดยกำหนดความถี่และเพดานจำนวนคำขอต่อเดือนเพื่อคุมค่าใช้จ่าย ทั้งสองแบบไม่เขียนทับฟิลด์ที่ผู้ดูแลระบบล็อกไว้ การซิงก์จึงเป็นแบบตามรอบเวลา มิใช่แบบเรียลไทม์ นอกจากนี้ ผู้ใช้แจ้งข้อมูลที่คลาดเคลื่อนและเสนอกิจกรรมใหม่ผ่านระบบ โดยผู้ดูแลระบบเป็นผู้ตรวจสอบก่อนเผยแพร่ และเมื่อแชตบอตค้นสถานที่จาก Google เพิ่มเติมเพราะฐานข้อมูลมีข้อมูลในพื้นที่ที่ผู้ใช้ถามไม่พอ ระบบบันทึกสถานที่เหล่านั้นไว้ในคิวให้ผู้ดูแลระบบพิจารณานำเข้า ข้อมูลจึงเพิ่มขึ้นตามสิ่งที่ผู้ใช้ถามถึงจริง
\item การเตรียมและจัดเก็บข้อมูล:
\begin{enumerate}
\item การทำความสะอาดและจัดรูปแบบข้อมูล (Data Preprocessing): ข้อมูลดิบที่รวบรวมมาจะถูกนำมาแปลงให้อยู่ในรูปแบบมาตรฐานของระบบ (Normalizing) เช่น การรวมประเภทสถานที่ย่อยจาก Google ให้อยู่ในหมวดหมู่ภาษาไทยที่ระบบใช้งาน (เช่น ร้านอาหาร คาเฟ่ วัด พิพิธภัณฑ์ สวนสาธารณะ ตลาด ที่พัก) และการตรวจสอบความสมบูรณ์ของข้อมูลก่อนบันทึกลงฐานข้อมูล
\item การแบ่งส่วนข้อมูลและการสร้างเวกเตอร์ (Chunking and Vector Embeddings): สำหรับการใช้งานร่วมกับ RAG ข้อมูลแต่ละรายการ (สถานที่ กิจกรรม และเนื้อหาในฐานความรู้) จะถูกแปลงเป็นเวกเตอร์ฝังตัว (Vector Embeddings) ด้วยโมเดล paraphrase-multilingual-MiniLM-L12-v2 (384 มิติ) ซึ่งประมวลผลบนเครื่องคอมพิวเตอร์ในเครื่อง (Local Inference) ผ่านไลบรารี sentence-transformers เมื่อมีการนำเข้าหรือแก้ไขข้อมูล ระบบจะล้างเวกเตอร์เดิมของรายการนั้น และคำนวณเวกเตอร์ใหม่เมื่อผู้ดูแลระบบสั่งสร้างดัชนีใหม่ (Reindex) จากหน้าจอผู้ดูแลระบบ
\item การจัดเก็บในฐานข้อมูลเวกเตอร์: เวกเตอร์ที่สร้างขึ้นจะถูกจัดเก็บไว้ในฐานข้อมูล PostgreSQL ตัวเดียวกับที่จัดเก็บข้อมูลหลักของระบบ โดยอาศัยส่วนขยาย pgvector ในการจัดเก็บและคำนวณค่าความคล้ายคลึงเชิงมุม (Cosine Similarity) การเลือกใช้ฐานข้อมูลเดียวกันแทนการติดตั้งฐานข้อมูลเวกเตอร์แยกต่างหาก ช่วยลดความซับซ้อนของระบบในระยะต้นแบบ และยังคงประสิทธิภาพการค้นคืนที่เพียงพอสำหรับขนาดข้อมูลของกรณีศึกษานี้
\end{enumerate}
\end{enumerate}

\subsection{วิเคราะห์และออกแบบระบบ}

การพัฒนาระบบแบ่งออกเป็น 3 องค์ประกอบที่พัฒนาและรันแยกจากกัน (Independent Services) ได้แก่ ส่วนเว็บแอปพลิเคชัน (Frontend) ส่วนบริการข้อมูลหลัก (Backend API) และส่วนบริการปัญญาประดิษฐ์ (AI Service) ซึ่งครอบคลุมทั้งระบบ แชตบอต และระบบวางแผนการเดินทาง

\begin{enumerate}
\item การออกแบบสถาปัตยกรรมระบบ: ระบบที่พัฒนาในงานวิจัยนี้ใช้แนวทางสถาปัตยกรรมไมโครเซอร์วิส (Microservices Architecture) ซึ่งประกอบด้วย 3 บริการหลักที่เป็นอิสระต่อกัน คือ Frontend Service, Backend API Service และ AI Service สื่อสารระหว่างกันผ่าน RESTful API เท่านั้น โดยรายละเอียดเทคโนโลยีและหน้าที่ของแต่ละบริการอธิบายไว้ในบทที่ 4 หัวข้อสถาปัตยกรรมของระบบ

\begin{figure}[H]
\centering
\adjustbox{max width=0.95\textwidth, max height=10cm}{%
\begin{tikzpicture}[
  box/.style={draw, thick, fill=white, align=center, font=\small, inner sep=5pt, minimum height=1.1cm},
  ext/.style={draw, thick, dashed, fill=white, align=center, font=\small, inner sep=4pt, minimum width=3.1cm, minimum height=0.9cm},
  flow/.style={thick, -{Stealth[length=2.4mm]}},
  lab/.style={font=\scriptsize, fill=white, inner sep=1.5pt, align=center}
]
\node[box, minimum width=3.4cm] (u1) at (-2.4,0) {นักท่องเที่ยว};
\node[box, minimum width=3.4cm] (u2) at (2.4,0) {ผู้ดูแลระบบ};
\node[box, minimum width=7.6cm] (front) at (0,-2.2) {Frontend\\เว็บแอปพลิเคชัน (React)};
\node[box, minimum width=3.8cm] (back) at (-3.2,-4.9) {Backend API\\(Node.js / Express)};
\node[box, minimum width=3.8cm] (ai) at (3.2,-4.9) {AI Service\\(Python / FastAPI)\\แชตบอตและระบบวางแผนทริป};
\node[cylinder, shape border rotate=90, aspect=0.25, draw, thick, fill=white, align=center, font=\small, minimum height=1.9cm, minimum width=3.6cm] (db) at (0,-7.9) {ฐานข้อมูล\\PostgreSQL + pgvector};

\node[ext] (gp) at (-7.7,-3.7) {Google Places API};
\node[ext] (cl) at (-7.7,-4.9) {Cloudinary\\(รูปภาพ)};
\node[ext] (em) at (-7.7,-6.1) {บริการอีเมล};
\node[ext] (llm) at (7.7,-4.9) {โมเดลภาษาขนาดใหญ่\\(LLM)};
\node[ext] (gp2) at (7.7,-3.4) {Google Places API\\(ค้นเพิ่มให้แชตบอต)};

\draw[flow] (u1.south) -- (u1 |- front.north);
\draw[flow] (u2.south) -- (u2 |- front.north);
\draw[flow] (front.south -| back) -- (back.north);
\draw[flow] (front.south -| ai) -- (ai.north);
\draw[flow] (back.east) -- node[lab, above] {ส่งต่อคำสั่ง\\ผู้ดูแลระบบ} (ai.west);
\draw[flow] (back.south) |- (db.west);
\draw[flow] (ai.south) |- (db.east);
\draw[flow] (back.west) -- (gp.east);
\draw[flow] (back.west) -- (cl.east);
\draw[flow] (back.west) -- (em.east);
\draw[flow] (ai.east) -- (llm.west);
\draw[flow] (ai.north east) -- (gp2.west);

\begin{scope}[on background layer]
  \node[draw, thick, dotted, rounded corners, inner sep=0.45cm, fit=(front)(back)(ai)(db), label={[font=\scriptsize]below:เครื่องแม่ข่ายเดียว (Docker Compose) ให้บริการภายใต้โดเมนเดียวผ่าน HTTPS}] {};
\end{scope}
\end{tikzpicture}%
}
\caption{แผนภาพสถาปัตยกรรมระบบของงานวิจัย}
\end{figure}

\item การออกแบบ UX/UI ของเว็บแอปพลิเคชัน: กำหนดโฟลว์การใช้งาน (User Flow) ตั้งแต่การเข้าสู่ระบบ การป้อนความต้องการท่องเที่ยว ไปจนถึงการรับแผนการเดินทางที่ระบบสร้างให้ ออกแบบหน้าจอสำคัญ เช่น หน้าแรก หน้ารายละเอียดสถานที่/กิจกรรม หน้าแบบฟอร์มวางแผนทริป และหน้าแสดงผลแผนการเดินทาง โดยยึดหลักการออกแบบที่มีผู้ใช้เป็นศูนย์กลาง
\end{enumerate}

\subsection{พัฒนาระบบ}

\begin{enumerate}
\needspace{5\baselineskip}
\item การพัฒนาระบบวางแผนการเดินทาง (Trip Planning): ระบบวางแผนการเดินทางใช้กลไก LLM ร่วมกับ RAG และชั้นการคำนวณเชิงกำหนดแน่นอน (Deterministic Computation Layer) ทำงานร่วมกันเป็นกลไกหลักในการสร้างแผนการเดินทาง โดยมีขั้นตอนดังนี้

\begin{figure}[H]
\centering
\adjustbox{max width=0.85\textwidth, max height=11cm}{%
\begin{tikzpicture}[
  stp/.style={draw, thick, rounded corners=3pt, fill=white, align=center, inner sep=5pt, text width=8.2cm},
  flow/.style={thick, -{Stealth[length=2.4mm]}},
  lab/.style={font=\small, fill=white, inner sep=1.5pt}
]
\node[stp] (s1) at (0,0) {รับและตรวจความต้องการของผู้ใช้\\(วันเดินทาง งบประมาณ ความสนใจ อาหาร จุดจบทริป)};
\node[stp] (s2) at (0,-2.0) {ค้นคืนสถานที่ผู้สมัครด้วย RAG แบบผสม};
\node[stp] (s3) at (0,-3.8) {LLM คัดเลือกสถานที่และจัดสรรลงในแต่ละวัน};
\node[stp] (s4) at (0,-5.9) {ชั้นการคำนวณเชิงกำหนดแน่นอนจัดลำดับการเดินทาง\\คำนวณตารางเวลา และตรวจสอบข้อจำกัด};
\node[draw, thick, diamond, aspect=2.6, fill=white, align=center, inner sep=1pt] (s5) at (0,-8.4) {LLM-as-a-Judge\\ผ่านหรือไม่};
\node[stp] (s6) at (0,-10.9) {แผนการเดินทางฉบับสมบูรณ์};
\draw[flow] (s1) -- (s2);
\draw[flow] (s2) -- (s3);
\draw[flow] (s3) -- (s4);
\draw[flow] (s4) -- (s5);
\draw[flow] (s5) -- node[lab, right] {ผ่าน หรือครบจำนวนรอบ} (s6);
\draw[flow] (s5.east) -- (6.0,-8.4) -- node[lab, sloped, above] {ไม่ผ่าน: คัดเลือกใหม่} (6.0,-3.8) -- (s3.east);
\end{tikzpicture}%
}
\caption{แผนภาพขั้นตอนการทำงานของระบบวางแผนการเดินทาง}
\end{figure}

ระบบประมวลผลตามลำดับ 4 ขั้นตอนหลัก ได้แก่ การทำความเข้าใจความต้องการของผู้ใช้ (Query Understanding) การค้นคืนสถานที่ผู้สมัครด้วยเทคนิค RAG แบบผสม (Candidate Retrieval) การให้ LLM คัดเลือกสถานที่และจัดสรรลงในแต่ละวันโดยอ้างอิงเฉพาะสถานที่ผู้สมัครที่ให้ไป (Prompt Construction \& Generation) และการจัดลำดับการเดินทาง คำนวณตารางเวลา และตรวจสอบข้อจำกัดด้วยชั้นการคำนวณเชิงกำหนดแน่นอนที่ทำงานอิสระจาก LLM (Guardrails / Constraint Verification) ทั้งนี้ LLM ไม่ได้เป็นผู้กำหนดลำดับหรือเวลาของแผน หน้าที่ดังกล่าวเป็นของชั้นการคำนวณเชิงกำหนดแน่นอนทั้งหมด หลังจากนั้นแผนที่จัดตารางเวลาแล้วจะถูกตรวจคุณภาพการคัดเลือกสถานที่ด้วยกลไก LLM-as-a-Judge หากไม่ผ่านระบบจะวนกลับไปให้ LLM คัดเลือกใหม่ภายในจำนวนรอบที่กำหนด

ในขั้นค้นคืน ระบบจัดอันดับสถานที่ผู้สมัครตามระดับความเกี่ยวข้องกับความสนใจก่อน แล้วจึงเรียงภายในระดับเดียวกันด้วยคะแนนความนิยม (คะแนนรีวิวที่ถ่วงด้วยจำนวนรีวิว) ร่วมกับระยะห่างจากที่พัก และกรองร้านอาหารและคาเฟ่ตามช่วงราคาต่อคนให้ตรงกับระดับงบประมาณ นอกจากเงื่อนไขพื้นฐาน ผู้ใช้ระบุประเภทอาหารที่อยากกินได้ ซึ่งระบบใช้เลือกร้านมื้อกลางวันและมื้อเย็น และเพิ่มมื้อเช้าให้เมื่อเลือก ``อาหารเช้า'' ระบุจุดจบของวันสุดท้ายได้ เช่น สนามบิน สถานีขนส่ง หรือสถานีรถไฟ พร้อมเวลาออกเดินทาง ซึ่ง Scheduler ใช้ตัดเวลาของวันสุดท้ายให้ทันเวลาและจัดสถานที่ให้อยู่ระหว่างทาง และเว้นจุดเริ่มต้นไว้ได้เมื่อยังไม่ได้จองที่พัก โดยระบบเริ่มจากใจกลางเมืองแล้วแนะนำย่านที่ควรพักจากตำแหน่งของสถานที่ในแผน

\item การพัฒนาระบบ แชตบอต และระบบแนะนำ (Chatbot and Recommendation System):

\begin{figure}[H]
\centering
\adjustbox{max width=0.9\textwidth, max height=12cm}{%
\begin{tikzpicture}[
  stp/.style={draw, thick, rounded corners=3pt, fill=white, align=center, inner sep=5pt, text width=7.4cm},
  side/.style={draw, thick, rounded corners=3pt, fill=white, align=center, inner sep=5pt, text width=3.6cm, font=\small},
  term/.style={draw, thick, rounded corners=8pt, fill=white, align=center, inner sep=5pt, minimum width=2.4cm},
  flow/.style={thick, -{Stealth[length=2.4mm]}},
  lab/.style={font=\small, fill=white, inner sep=1.5pt}
]
\node[term] (s0) at (0,0) {ผู้ใช้ส่งคำถาม};
\node[stp] (s1) at (0,-1.7) {รับคำถามพร้อมประวัติการสนทนา\\(เก็บไว้ 4 ข้อความล่าสุด)};
\node[stp] (s2) at (0,-3.6) {LLM เขียนคำถามใหม่ให้สมบูรณ์ในตัวเอง\\(Query Rewriting)};
\node[stp] (s3) at (0,-5.7) {ค้นคืนข้อมูลแบบผสม (Hybrid Search)\\จากตาราง places, knowledge\_base และ events\\ตารางละ 3 รายการ};
\node[stp] (s4) at (0,-8.0) {LLM สร้างคำตอบจากข้อมูลที่ค้นคืนได้\\ส่งกลับแบบ Streaming};
\node[draw, thick, diamond, aspect=2.6, fill=white, align=center, inner sep=1pt] (d1) at (0,-10.6) {มีข้อมูล\\รองรับคำตอบหรือไม่};
\node[stp, text width=3.8cm, font=\small] (no) at (-4.3,-13.1) {ตอบด้วยข้อความมาตรฐาน\\ว่าไม่พบข้อมูล};
\node[stp, text width=4.4cm, font=\small] (yes) at (4.3,-13.1) {แสดงคำตอบพร้อมการ์ด\\สถานที่หรือกิจกรรมที่อ้างถึง};
\node[term] (end) at (0,-15.1) {จบ};
\node[side] (fail) at (8.2,-8.0) {แสดงข้อความขออภัย\\(LLM หรือฐานข้อมูลล้มเหลว)};
\node[side, dashed] (gfb) at (8.2,-4.4) {ถามแบบ ``แถว X'' แต่พบไม่ถึง 3 แห่ง: ค้นเพิ่มจาก Google Places};

\draw[flow] (s0) -- (s1);
\draw[flow] (s1) -- (s2);
\draw[flow] (s2) -- (s3);
\draw[flow] (s3) -- (s4);
\draw[flow] (s4) -- (d1);
\draw[flow] (d1.west) -| node[lab, pos=0.25, above] {ไม่มี} (no.north);
\draw[flow] (d1.east) -| node[lab, pos=0.25, above] {มี} (yes.north);
\draw[flow] (no.south) |- (end.west);
\draw[flow] (yes.south) |- (end.east);
\draw[flow] (s4.east) -- node[lab, above] {ล้มเหลว} (fail.west);
\draw[thick, dashed, {Stealth[length=2.4mm]}-{Stealth[length=2.4mm]}] (s3.east) -| node[lab, pos=0.3, below] {ค้นเพิ่ม และรวมผล} (gfb.south);
\draw[flow] (fail.south) -- (8.2,-16.5) -- (0,-16.5) -- (end.south);
\end{tikzpicture}%
}
\caption{แผนภาพขั้นตอนการทำงานของระบบ แชตบอต และระบบแนะนำ}
\end{figure}

ระบบประมวลผลคำถามของผู้ใช้ผ่าน 3 ขั้นตอนหลัก ได้แก่ การเขียนคำถามใหม่จากประวัติการสนทนาและค้นคืนข้อมูลที่เกี่ยวข้องจากสถานที่ ฐานความรู้ และกิจกรรมด้วยเทคนิค RAG แบบผสม (Information Retrieval) การให้ LLM สังเคราะห์คำตอบแบบทยอยแสดงผลพร้อมกลไกป้องกันการตอบข้อมูลที่ไม่มีหลักฐานรองรับ (Response Generation with Hallucination Guard) และการแสดงการ์ดสถานที่หรือกิจกรรมแนะนำจากผลการค้นคืน เฉพาะแหล่งข้อมูลที่คำตอบอ้างถึง (Recommendation)

สำหรับคำถามที่ระบุพื้นที่ เช่น ``ร้านอาหารแถวบึงแก่นนคร'' หรือ ``ฉันพักอยู่แถว มข.'' ระบบแยกชื่อพื้นที่ออกจากคำถาม แล้วหาตำแหน่งของพื้นที่นั้นจากตำบล อำเภอ หรือสถานที่สำคัญในฐานข้อมูล และจัดอันดับสถานที่รอบตำแหน่งดังกล่าวตามระยะทาง แทนการเทียบเพียงชื่อและแท็กของสถานที่ หากฐานข้อมูลมีสถานที่ในพื้นที่นั้นไม่ถึง 3 แห่ง ระบบค้นเพิ่มจาก Google Places API และติดป้ายว่าเป็นข้อมูลที่ยังไม่ผ่านการตรวจสอบของทีมงาน โดยจำกัดจำนวนการค้นต่อวันและเก็บผลไว้ 1 ชั่วโมงเพื่อคุมค่าใช้จ่าย สถานที่จาก Google ที่แสดงไปถูกบันทึกเป็นรายการให้ผู้ดูแลระบบพิจารณานำเข้า นอกจากนี้ระบบตอบคำถามเรื่องเวลาเปิด-ปิดได้ถูกต้องขึ้น รวมถึงร้านที่เปิดข้ามคืน และกรองร้านมังสวิรัติหรือวีแกนจากข้อมูลสิ่งอำนวยความสะดวกของสถานที่

\item การพัฒนา Frontend: พัฒนาด้วยเฟรมเวิร์ก React ร่วมกับเครื่องมือ Vite เพื่อสร้างส่วนติดต่อผู้ใช้ที่โต้ตอบได้และรองรับการใช้งานข้ามอุปกรณ์ (Responsive Web Design) ทั้งบน Desktop, Tablet และ Mobile
\item การพัฒนา Backend และการเชื่อมต่อกับบริการต่าง ๆ: ออกแบบ RESTful API สำหรับดึงข้อมูลจากบริการวางแผนการเดินทาง ระบบแนะนำ และข้อมูลกิจกรรม/สถานที่ พัฒนาหน้าจอสำหรับผู้ดูแลระบบ (Admin Dashboard) เพื่อจัดการข้อมูลสถานที่ กิจกรรม ฐานความรู้ และรหัส QR สำหรับระบบสะสมคะแนน รวมถึงระบบบัญชีผู้ใช้ที่ยืนยันตัวตนด้วยเซสชันฝั่งเซิร์ฟเวอร์ คิวตรวจสอบรายงานและคำขอจากผู้ใช้ และเครื่องมือบริหารบัญชีและพอยท์ ส่วนการดูแลข้อมูลสถานที่ประกอบด้วยการนำเข้าจาก Google ตามชื่อหรือตามพื้นที่ การซิงก์ทั้งแบบสั่งเองและแบบอัตโนมัติตามรอบเวลาซึ่งทำงานเบื้องหลังใน Backend และคิวสถานที่ที่ผู้ใช้ถามถึงผ่านแชตบอต โดยปุ่มที่เรียก API ที่มีค่าบริการ เช่น ซิงก์ นำเข้า และสร้างคำอธิบายด้วย AI แสดงค่าใช้จ่ายโดยประมาณต่อครั้งให้ผู้ดูแลระบบเห็นก่อนกด
\item การทดสอบด้านการใช้งาน (Usability Testing): ทดสอบกับกลุ่มผู้ใช้เป้าหมาย (นักศึกษา/นักท่องเที่ยวในจังหวัดขอนแก่น) โดยใช้เกณฑ์ประเมินด้านความง่ายในการใช้งาน ความพึงพอใจและความเข้าใจในการนำเสนอข้อมูล
\end{enumerate}

\subsection{ทดสอบและประเมินผล}

\begin{enumerate}
\item การทดสอบระบบวางแผนการเดินทาง: ทดสอบ 2 ระดับ ระดับแรกเป็นชุดทดสอบอัตโนมัติ (Unit Test) ของ Judge, Trip Planner และ Scheduler ที่จำลองคำตอบของ LLM (Mock) เพื่อตรวจว่าการคัดเลือกสถานที่ การจัดตารางเวลา และการตัดสินผ่านของ Judge ทำงานตามที่ออกแบบ ระดับที่สองเป็นการทดลองแบบออฟไลน์ที่เรียกโมเดลจริงเพื่อเปรียบเทียบ Prompt ของตัวสร้างแผน ด้วยกรณีทดสอบที่กำหนดตายตัว และวัดผลด้วยตัวชี้วัด Pass Rate ที่ปรับจากกรอบแนวคิด TravelPlanner \parencite{xie2024} 3 ตัว คือ Environment (แผนไม่มีรหัสสถานที่ที่ไม่อยู่ในฐานข้อมูลหรือซ้ำกัน ซึ่งใช้ยืนยันกลไกป้องกัน Hallucination) Hard (สถานที่ที่ผู้ใช้กำหนดว่าต้องไปอยู่ครบในแผนสุดท้าย) และ Commonsense (ทุกวันมีมื้ออาหาร และไม่เลือกสถานที่ที่ปิดในวันนั้น) ร่วมกับคะแนน 1--5 จาก LLM ผู้ประเมิน 3 ตัวที่แยกจาก Judge ในระบบจริง ใน 5 ด้าน ได้แก่ ความตรงกับความสนใจ ความตรงกับเงื่อนไข ความหลากหลาย ความสอดคล้องกันภายในแต่ละวัน และคุณภาพของการเลือกสถานที่ แล้วทดสอบสมมติฐานด้วยการวิเคราะห์ความแปรปรวนแบบวัดซ้ำสองทาง (Two-way Repeated-Measures ANOVA) ใน 3 ประเด็น คือ Prompt ที่ต่างกันส่งผลต่อคุณภาพของแผนหรือไม่ ผู้ประเมินแต่ละตัวให้คะแนนเข้มงวดต่างกันหรือไม่ และผู้ประเมินเห็นผลของ Prompt ตรงกันหรือไม่

\item การทดสอบแชตบอต: ใช้ชุดคำถามมาตรฐาน (Golden Set) ที่เรียกโมเดลจริงเพื่อตรวจสอบเชิงโครงสร้างของคำตอบ ร่วมกับชุดทดสอบอัตโนมัติที่จำลองคำตอบของ LLM และของ Google Places API เพื่อตรวจการค้นตามพื้นที่ การค้นเพิ่มจาก Google และการตอบเรื่องเวลาเปิด-ปิด และผลประเมินความพึงพอใจของผู้ใช้ด้านแชตบอตจากแบบสอบถาม

\item การทดสอบส่วนเว็บ: นอกจากกรณีทดสอบเชิงหน้าที่ที่กำหนดไว้ล่วงหน้า Backend มีชุดทดสอบอัตโนมัติที่เรียกเส้นทาง API ทั้งแบบสาธารณะและแบบที่ต้องเข้าสู่ระบบด้วยฐานข้อมูลจำลอง และตรวจว่าทุกเส้นทางของ Backend มีอยู่ในเอกสาร OpenAPI ครบถ้วน

\item การประเมินความพึงพอใจของผู้ใช้: ใช้แบบสอบถามในระบบที่ผู้ใช้ตอบหลังทดลองใช้งาน ใช้มาตรประมาณค่า 5 ระดับ จำนวน 16 ข้อใน 6 ด้าน ได้แก่ ด้านการทำงานของระบบ ด้านแชตบอต AI ด้านการวางแผนเที่ยวด้วย AI ด้านความง่ายในการใช้งาน ด้านการออกแบบและประสิทธิภาพ และด้านความพึงพอใจโดยรวม พร้อมคำถามปลายเปิด และแปลผลค่าเฉลี่ยด้วยเกณฑ์ช่วงกว้าง 0.80 \parencite{best1981} รายละเอียดของเครื่องมือและเกณฑ์อยู่ในหัวข้อการทดสอบระบบของบทที่ 4
\end{enumerate}

\subsection{ปรับปรุงและสรุปผล}

\begin{enumerate}
\item การวิเคราะห์ผลลัพธ์ของโมเดล: เปรียบเทียบผลการประเมินที่ได้จากหัวข้อการทดสอบและประเมินผล โดยเน้นการวิเคราะห์เชิงลึก 2 ส่วน คือ Trip Planning วิเคราะห์ว่าการผสานกลไก LLM เข้ากับชั้นการคำนวณเชิงกำหนดแน่นอน (Deterministic Backstop) ภายใต้สถาปัตยกรรมแบบแยกส่วนบริการสามารถยกระดับคุณภาพของแผนการเดินทาง (Itinerary Quality) ได้อย่างมีประสิทธิภาพตามเป้าหมายที่ตั้งไว้หรือไม่ เมื่อเทียบกับแนวทางที่ใช้ LLM สร้างแผนโดยไม่มีชั้นตรวจสอบเพิ่มเติม; และ แชตบอต วิเคราะห์ว่ากลไก RAG ร่วมกับการค้นคืนแบบผสม (Hybrid Search) ช่วยให้ระบบให้คำแนะนำที่ถูกต้องและมีความเกี่ยวข้องของคำตอบสูงขึ้นโดยไม่เพิ่มอัตราการเกิด Hallucination หรือไม่ เมื่อเทียบกับแนวทางค้นคืนด้วยความคล้ายคลึงเชิงเวกเตอร์เพียงอย่างเดียว

\item การวิเคราะห์เชิงพฤติกรรมและการยอมรับของผู้ใช้: วิเคราะห์ข้อมูลความพึงพอใจของผู้ใช้ โดยพิจารณาว่าผู้ใช้รับรู้ถึงประโยชน์หลักของการใช้ AI อย่างไร (เช่น การเข้าถึงข้อมูลที่รวดเร็วขึ้น การได้รับแผนการเดินทางที่ตรงความต้องการโดยไม่ต้องค้นหาข้อมูลเอง เวลาที่ใช้ในการวางแผนสั้นลง) และระบุข้อกังวลหลักของผู้ใช้ (เช่น ความถูกต้องของข้อมูลที่ AI ให้ ความเป็นส่วนตัวและความปลอดภัยของข้อมูล การพึ่งพาเทคโนโลยีสูง) ซึ่งข้อมูลเหล่านี้สามารถนำไปสู่การปรับปรุงความโปร่งใสของระบบ (เช่น การแสดงแหล่งที่มาของคำแนะนำ) และการสร้างความเชื่อมั่นในระบบ

\item การสรุปผลและข้อเสนอแนะ: สรุปผลลัพธ์หลักที่แสดงให้เห็นถึงความสำเร็จของการบูรณาการ LLM และเทคนิค RAG เข้ากับสถาปัตยกรรมแบบแยกส่วนบริการในการวางแผนการเดินทางและการแนะนำการท่องเที่ยว พร้อมระบุข้อจำกัดที่พบในงานวิจัย ได้แก่ (1) ปัญหา Cold-start สำหรับผู้ใช้ใหม่ที่ยังไม่มีประวัติความสนใจ (2) ขนาดชุดข้อมูลที่จำกัดเฉพาะพื้นที่จังหวัดขอนแก่น (3) มาตรการด้านความปลอดภัยของเว็บแอปพลิเคชันที่ยังไม่ครบถ้วน เช่น นโยบาย Content Security Policy และการล็อกบัญชีหลังเข้าสู่ระบบผิดซ้ำ จากข้อจำกัดดังกล่าว งานวิจัยนี้เสนอแนวทางสำหรับการพัฒนาต่อยอดในอนาคต ได้แก่ การขยายขอบเขตให้ครอบคลุมตัวชี้วัดด้านความยั่งยืนเชิงปริมาณที่ยังไม่ได้นำมาใช้ในต้นแบบปัจจุบัน (เช่น Carbon Footprint จากการเดินทาง และผลกระทบทางเศรษฐกิจต่อผู้ประกอบการท้องถิ่น) การเสริมสถาปัตยกรรมด้วย API Gateway และ Message Broker เพื่อรองรับการขยายตัวของระบบ การเสริมมาตรการความปลอดภัยของเว็บแอปพลิเคชัน และการพัฒนาไปสู่ระบบแนะนำแบบสนทนาที่ซับซ้อนยิ่งขึ้น (Conversational Recommender System) ที่สามารถปรับแผนการเดินทางแบบโต้ตอบกับผู้ใช้ได้แบบต่อเนื่อง
\end{enumerate}

% -----------------------------------------------------------
\section{เครื่องมือที่ใช้ในการดำเนินงาน}
% -----------------------------------------------------------


\subsection{ฮาร์ดแวร์}

% สเปกจากหน้าแพ็กเกจ VPS-1 ของ OVHcloud ที่ผู้จัดทำเช่าใช้งานจริง
ระบบที่ให้บริการจริงทำงานบนเซิร์ฟเวอร์เสมือน (Virtual Private Server: VPS) แพ็กเกจ VPS-1 ของ OVHcloud เพียงเครื่องเดียว ซึ่งรันบริการ Frontend, Backend และ AI Service รวมถึงการสร้างเวกเตอร์ฝังตัวแบบประมวลผลในเครื่อง ผ่าน Docker Compose โดยมี Caddy ทำหน้าที่ส่งต่อคำขอ ส่วนฐานข้อมูลและรูปภาพใช้บริการภายนอก (Supabase และ Cloudinary) จึงไม่ได้ใช้ทรัพยากรของเครื่องนี้ รายละเอียดของเครื่องแสดงดังตารางที่ \ref{tab:hardware}

\begin{table}[H]
\centering
\caption{รายละเอียดเซิร์ฟเวอร์ที่ใช้ให้บริการระบบ}
\label{tab:hardware}
\begin{tabular}{|p{4.5cm}|p{8.5cm}|}
\hline
\textbf{รายการ} & \textbf{รายละเอียด} \\
\hline
ผู้ให้บริการและแพ็กเกจ & OVHcloud VPS-1 \\ \hline
หน่วยประมวลผล & 2 vCores \\ \hline
หน่วยความจำ & 4 GB RAM \\ \hline
พื้นที่จัดเก็บ & 40 GB SSD NVMe พร้อมสำรองข้อมูลรายวัน (เก็บย้อนหลัง 24 ชั่วโมง) \\ \hline
เครือข่าย & แบนด์วิดท์สาธารณะ 500 Mbps ไม่จำกัดปริมาณการรับส่งข้อมูล \\ \hline
ค่าบริการ & 180 บาทต่อเดือน \\ \hline
\end{tabular}
\end{table}

% กันหัวข้อและประโยคนำไม่ให้ค้างท้ายหน้า เมื่อตารางที่ตามมาลงหน้าเดิมไม่พอ
\subsection{ซอฟต์แวร์}

เครื่องมือด้านซอฟต์แวร์ที่ใช้ในการพัฒนาโครงงานแสดงดังตารางที่ \ref{tab:software}

\begin{table}[H]
\centering
\caption{รายการซอฟต์แวร์ที่ใช้ในการพัฒนา}
\label{tab:software}
\begin{tabular}{|c|p{4.3cm}|p{6.5cm}|}
\hline
\textbf{ลำดับ} & \textbf{ซอฟต์แวร์} & \textbf{วัตถุประสงค์} \\
\hline
1 & React (Vite) & พัฒนา Frontend Service (Web) \\ \hline
2 & Node.js (Express) & พัฒนา Backend API Service หลัก \\ \hline
3 & Python (FastAPI) & พัฒนา AI Service สำหรับระบบ แชตบอต และระบบวางแผนการเดินทาง (ใช้ RAG) \\ \hline
4 & PostgreSQL (Supabase) + pgvector & จัดเก็บข้อมูลหลักของระบบและเวกเตอร์ฝังตัว (Vector Embeddings) สำหรับ RAG ในฐานข้อมูลเดียวกัน \\ \hline
5 & sentence-transformers (paraphrase-multilingual-MiniLM-L12-v2) & สร้างเวกเตอร์ฝังตัวแบบประมวลผลในเครื่อง (Local Inference) \\ \hline
6 & gemini-3.8-flash (ผ่าน Vertex AI ของ Google Cloud) & โมเดลภาษาขนาดใหญ่ (LLM) สำหรับแชตบอต การวางแผนทริป การตรวจคุณภาพแผน (Judge) การดึงข้อมูลกิจกรรมจากโพสต์ และการร่างคำอธิบายสถานที่ \\ \hline
7 & Google Places API & ดึงข้อมูลตั้งต้น (Data Seeding) ของสถานที่ท่องเที่ยว \\ \hline
8 & Cloudinary & จัดเก็บและส่งมอบรูปภาพของสถานที่ กิจกรรม และของรางวัล \\ \hline
9 & Docker Compose และ Caddy & รวมบริการทั้งสามส่วนบนเซิร์ฟเวอร์เครื่องเดียว พร้อมใบรับรอง HTTPS และการส่งต่อคำขอ \\ \hline
10 & GitHub Actions & ตรวจสอบการสร้างและทดสอบโค้ด และสั่งปรับปรุงระบบบนเซิร์ฟเวอร์อัตโนมัติ \\ \hline
11 & claude-sonnet-5.5, gpt-5.6-terra และ deepseek-v4-pro & LLM ผู้ประเมินอิสระในการทดลองเปรียบเทียบ Prompt ของระบบวางแผนการเดินทางแบบออฟไลน์ โดย claude-sonnet-5.5 เป็นผู้ประเมินหลัก ส่วนอีกสองตัวใช้วิเคราะห์ความไวของผล (ไม่ได้ใช้ในระบบที่ให้บริการจริง) \\
\hline
\end{tabular}
\end{table}

\subsection{ค่าใช้จ่ายในการดำเนินงาน}

% ตัวเลขจากผู้จัดทำ: VM ของ OVHcloud เดือนละ 180 บาท (นับตั้งแต่ Deploy ต้นเดือนกันยายน 2569 ตามบทที่ 5),
% โดเมนจาก Cloudflare 200 บาท, ค่า Vertex AI จากหน้า Billing ของ Google Cloud ช่วง 1--15 ต.ค. 2569 (เริ่มมีค่าใช้จ่ายวันที่ 10)
% TODO: Google Places API -- เครดิตทดลองใช้ถูกใช้ไปแล้ว 2,182.99 บาท (10,066.05 - 7,883.06) แต่เดือนตุลาคมมีเพียง Vertex AI 135.24 บาท
%       ส่วนที่เหลือน่าจะมาจากการดึงข้อมูล Google Places ก่อนเดือนตุลาคม -- ตรวจยอดจริงใน Billing แล้วเพิ่มแถว
ค่าใช้จ่ายในการพัฒนาและให้บริการระบบ นับถึงวันที่ 15 ตุลาคม 2569 แสดงดังตารางที่ \ref{tab:cost} บริการที่ใช้แพ็กเกจฟรีได้แก่ ฐานข้อมูล Supabase บริการจัดเก็บรูปภาพ Cloudinary และ Serper API ส่วนค่าใช้ Vertex AI ชำระด้วยเครดิตทดลองใช้ฟรี (Free Trial Credits) ของ Google Cloud จึงไม่ได้ชำระเป็นเงินจริง เงินที่คณะผู้จัดทำชำระจริงจึงมีเพียงค่าเซิร์ฟเวอร์และค่าโดเมน รวม 560 บาท

\begin{table}[H]
\centering
\caption{ค่าใช้จ่ายในการดำเนินงาน (ข้อมูลถึงวันที่ 15 ตุลาคม 2569)}
\label{tab:cost}
\begin{tabular}{|p{4.6cm}|p{2.4cm}|p{4.2cm}|r|}
\hline
\textbf{รายการ} & \textbf{ผู้ให้บริการ} & \textbf{ช่วงเวลา} & \textbf{บาท} \\
\hline
เซิร์ฟเวอร์เสมือน (VPS) สำหรับ Frontend, Backend และ AI Service & OVHcloud & ก.ย.--ต.ค. 2569 (2 เดือน เดือนละ 180 บาท) & 360.00 \\ \hline
โดเมน dinokhonkaen.app & Cloudflare & ค่าจดทะเบียนโดเมน & 200.00 \\ \hline
LLM gemini-3.8-flash ผ่าน Vertex AI & Google Cloud & ต.ค.2569 (ชำระด้วยเครดิตทดลองใช้ฟรี) & 139.00 \\ \hline
ฐานข้อมูล (Supabase), รูปภาพ (Cloudinary) และ Serper API & -- & แพ็กเกจฟรี & 0.00 \\ \hline
\multicolumn{3}{|l|}{\textbf{รวมค่าใช้จ่ายทั้งหมด}} & \textbf{699.00} \\ \hline
\multicolumn{3}{|l|}{\textbf{ชำระเป็นเงินจริง (ไม่รวมส่วนที่ใช้เครดิต)}} & \textbf{560.00} \\ \hline
\end{tabular}
\end{table}

% -----------------------------------------------------------
\section{แผนการดำเนินงาน}
% -----------------------------------------------------------

แผนการดำเนินงานตลอดระยะเวลาโครงงานแสดงดังตารางที่ \ref{tab:gantt} โดยช่องที่แรเงาหมายถึงเดือนที่ดำเนินงานในขั้นตอนนั้น

\begin{table}[H]
\centering
\caption{แผนการดำเนินงาน (Gantt Chart)}
\label{tab:gantt}
\begin{tabular}{|p{5.5cm}|c|c|c|c|c|}
\hline
 & \multicolumn{5}{c|}{ปี 2569} \\
\hline
การดำเนินงาน & 6 & 7 & 8 & 9 & 10 \\
\hline
1. การศึกษาวรรณกรรมและทฤษฎีที่เกี่ยวข้อง
 & \cellcolor{gray!60} & & & & \\
\hline
2. การเก็บรวบรวมและเตรียมข้อมูล
 & \cellcolor{gray!60} & \cellcolor{gray!60} & & & \\
\hline
3. การออกแบบและพัฒนาระบบ
 & & \cellcolor{gray!60} & \cellcolor{gray!60} & & \\
\hline
4. การทดสอบและประเมินผลระบบ
 & & & \cellcolor{gray!60} & \cellcolor{gray!60} & \cellcolor{gray!60} \\
\hline
5. การวิเคราะห์และสรุปผล
 & & & & & \cellcolor{gray!60} \\
\hline
\end{tabular}
\end{table}

จากตารางที่ \ref{tab:gantt} การดำเนินงานเริ่มจากการศึกษาวรรณกรรมและเก็บรวบรวมข้อมูลในเดือนมิถุนายน--กรกฎาคม 2569 ตามด้วยการออกแบบและพัฒนาระบบในเดือนกรกฎาคม--สิงหาคม การทดสอบและประเมินผลระบบในเดือนสิงหาคม--ตุลาคม และปิดท้ายด้วยการวิเคราะห์และสรุปผลในเดือนตุลาคม 2569